\scanpage{73}
Considérons la forme quadratique $\unsure{q_0|q}$ sur $\widetilde{\mathfrak g}$
\[
\left\{\begin{aligned}
q&=Q\mid\widetilde{\mathfrak g}:\widetilde{\mathfrak g}\longrightarrow\mathcal O_S\\
q(u)&=-\det u
\end{aligned}\right.\tag*{(43)}
\]
alors \deleted{$\unsure{q}$} $\varphi$ est donc associée à la forme $q$
\[
\left\{\begin{aligned}
q(x+y)&=q(x)\unsure{{}\,|+}q(y)+\varphi(x,y),\\
\varphi(u,u)&=2q(u)
\end{aligned}\right.\tag*{(44)}
\]
d'ailleurs $q$ est lisse -- $\varphi$ est dualisante sauf en car. $2$.

[La relation avec Killing est
\[
\left\{\begin{aligned}
\operatorname{Tr}(\operatorname{ad}_{\widetilde{\mathfrak g}}(u)\operatorname{ad}_{\widetilde{\mathfrak g}}(v))&=4\varphi(u,v)\\
\operatorname{Tr}(\operatorname{ad}_{\widetilde{\mathfrak g}}(u)^2)&=8q(u)
\end{aligned}\right.\tag*{(45)}
\]
pour $u,v$ sections de $\widetilde{\mathfrak g}$. \enspace]

On définit aussi une base
\[
\omega\in\Gamma\underline{\operatorname{\text{dét}}}\,\widetilde{\mathfrak g}\simeq\Gamma\underline{\operatorname{\text{dét}}}\,M\tag*{(46)}
\]
(le dernier iso provient de la suite exacte (35$'$) \unsure{en passant au dét}) compte tenu de l'isom local
\[
M\simeq\underline{\operatorname{End}}(V)\simeq\check V\otimes V
\]
et de l'isomorphisme canonique
\[
\begin{aligned}
\det(\check V\otimes V)&\simeq\det\check V\otimes\det V\simeq(\underline{\det}V)^{-1}\otimes\underline{\det}V\\
&\simeq\mathcal O_S
\end{aligned}\tag*{(47)}
\]
% archive sheet number: 72
